% Mathematical construction and formalization targets. This file does not
% register any new checked Lean theorem or alter a numerical Lean evaluator.
\section{Continue the logarithm from one to $i$}
\label{sec:logarithm-continuation}

The exponential connects real growth with rotation. The logarithm makes the
reverse connection: continuing its local computation to the rational complex
point $i$ supplies another computation of the geometric constant $\pi$.
The continuation algorithm must not use that constant to choose its values.

\subsection{A local computation and its tail}
For a rational complex input $w$, start with
\[
 \ell_N(w)=\sum_{r=1}^N\frac{(-1)^{r+1}w^r}{r}.
\]
If $|w|\le q<1$, with a supplied rational $q$, the omitted tail is at most
\[
 \frac{q^{N+1}}{(N+1)(1-q)}.
\]
This gives rational boxes for the local function $\ell(w)$. Its agreement
with $L(1+w)$ for real $-1<w<1$ follows by integrating the finite geometric
identity for $1/(1+w)$ and retaining its remainder. Thus the series is a
new representation of the existing logarithm, not a replacement for its
reciprocal-integral definition.

At a nonzero rational complex centre $z$, an increment $h$ satisfying
$|h/z|<1$ changes the continued logarithm by $\ell(h/z)$. The local
identity $\ell'(w)=1/(1+w)$, with its uniform finite remainder bound,
justifies this update. Starting from the value zero at one, a finite chain
of such steps computes a specified continuation. The path is part of the
construction; one must not silently identify continuations winding differently
around zero.

\subsection{Four explicit steps, with no trigonometric input}
Use the straight segment $\gamma(t)=1-t+it$, $0\le t\le1$, subdivided at
four equal rational steps:
\[
 z_0=1,\quad z_1=\frac{3+i}{4},\quad z_2=\frac{1+i}{2},\quad
 z_3=\frac{1+3i}{4},\quad z_4=i.
\]
For $w_j=(z_{j+1}-z_j)/z_j$, the four local inputs are
\[
 w_0=\frac{-1+i}{4},\qquad
 w_1=\frac{-1+2i}{5},\qquad
 w_2=\frac{i}{2},\qquad
 w_3=\frac{1+2i}{5}.
\]
Their squared moduli are $1/8,1/5,1/4,1/5$. In particular every step admits
the same rational bound $q=1/2$. The path stays apart from zero because
$|\gamma(t)|^2=(1-t)^2+t^2\ge1/2$.

Define the endpoint continuation by summing the four local increments:
\[
 \mathcal L_\gamma(i)=\sum_{j=0}^3\ell(w_j),\qquad
 \pi_{\log}=2\operatorname{Im}\mathcal L_\gamma(i).
\]
This is a definition of an independent computation, before identifying its
value with geometric $\pi$. Its finite approximation and error radius are
\[
 p_N=2\operatorname{Im}\sum_{j=0}^3\ell_N(w_j),\qquad
 R_N=\frac{8}{(N+1)2^N}.
\]
The rational interval $[p_N-R_N,p_N+R_N]$ encloses $\pi_{\log}$.
Taking the finite intersections of these intervals for $1\le n\le N$
gives nested outputs with widths at most $2R_N$. Compatibility and shrinking
follow from the geometric tail estimates themselves, without using geometric
$\pi$ as an anchor or assuming the desired identification.

Everything in these finite instructions is rational complex arithmetic.
There is no call to arctangent, trigonometry, a pre-existing complex logarithm,
or an evaluator for $\pi$. Evaluating all four prefixes gives the following
outward-rounded rational enclosures (the numerical calculation is not yet
an extracted Lean program):
\[
\begin{array}{c|c}
 N&\hbox{enclosure of }\pi_{\log}\\\hline
 8 &[3.137914622420,\ 3.144859066866]\\
 16&[3.141584712311,\ 3.141599073525]\\
 32&[3.141592653527,\ 3.141592653641]\\
 64&[3.141592653589,\ 3.141592653590]
\end{array}
\]
The fixed decimal display eventually hides further improvement of the exact
rational bounds; it is not the precision used to generate them.

\subsection{Identify the result with geometric area}
Each local series is compared with the integral of the reciprocal along
its segment by the finite geometric remainder identity. Adding the four
comparisons gives
\[
 \mathcal L_\gamma(i)
 =\int_0^1\frac{-1+i}{1-t+it}\,dt
 =\int_0^1\frac{2t-1+i}{(1-t)^2+t^2}\,dt.
\]
The complex expression means two real computations, not an assumed general
contour-integral operator. Their rational denominators are bounded below
by $1/2$. The real integrand is increasing, while the imaginary integrand
is increasing up to $1/2$ and decreasing thereafter, so the existing monotone
and finite-piece constructions are appropriate.

The real part cancels under $t\mapsto1-t$. In the imaginary part the rational
linear substitution $u=2t-1$ gives
\[
 \operatorname{Im}\mathcal L_\gamma(i)
 =\int_{-1}^{1}\frac{du}{1+u^2}
 =2\int_0^1\frac{du}{1+u^2}
 \simeq2A(1).
\]
Consequently the intended equivalence is
\[
 \boxed{\mathcal L_\gamma(i)\simeq 2iA(1)=\frac{i\pi}{2},\qquad
        \pi_{\log}\simeq4A(1)=\pi.}
\]
The arctangent computation is used in this identification proof, not by the
new evaluator. No general homotopy or winding-number theorem is required
for this particular path. Agreement after a finite subdivision is a useful
first continuation theorem; comparison of paths with different winding belongs
to the later complex-path chapter.

\subsection{The exponential and the circle are part of the same story}
The factorial exponential and the continued logarithm should be compared
locally before invoking geometric $\pi$. The local inverse identity and
finite multiplication give
\[
 E(\ell(w))\simeq1+w,\qquad
 E(\mathcal L_\gamma(i))\simeq\prod_{j=0}^3(1+w_j)=i.
\]
The final product telescopes because $1+w_j=z_{j+1}/z_j$. It is exact rational
complex algebra. Combined with the geometric identification, this yields
$E(i\pi/2)\simeq i$ without defining either side through the other.

The fuller trigonometric connection is the equality, on the existing chart,
\[
 E(i\pi x)\simeq C(x)+iS(x).
\]
Both coordinates must be identified with the original inverse-arctangent
constructions, not just assigned their names. One useful proof compares the
factorial exponential with the rational circle path
\[
 P(u)=\frac{1-u^2+2iu}{1+u^2}.
\]
Both $E(2iA(u))$ and $P(u)$ have initial value one and satisfy
$Y'=2iY/(1+u^2)$. The quantitative finite uniqueness argument already described
above identifies them. Alternatively compare their normalized rotation
systems on a bounded rational interval. In either route, composition at the
computed input $A(u)$ uses its quantitative continuity bounds.

This makes real growth, complex rotation and phase computations aspects of
one computational calculus. The same continuation interface can later support
complex powers $E(\alpha\mathcal L_\gamma(z))$, with the path and separation
conditions retained as explicit input data.

\subsection{Formalization and comparison targets}
The new continuation construction above is a mathematical specification and
an exact-rational numerical example. It is not yet registered as a checked
Lean development. The required native results are local series validity,
finite continuation and subdivision agreement, local exponential inversion,
and identification with the rational path integral and geometric arctangent.
They should reuse the existing logarithm, exponential, complex-box and
finite-integral machinery rather than create parallel foundations.

The paired endpoint target is the equivalence of the same $\pi_{\log}$
program with the original geometric $\pi$. The native proof follows the
finite reciprocal-integral calculation. The optional Mathlib proof must
first identify the local series and its branch continuation with Mathlib's
logarithm, then use its value at $i$. Calling its ready-made logarithm is
not an implementation of our continuation algorithm.

The bundled proof map should therefore distinguish the rational local-series
and continuation definitions, the exponential inverse and trigonometric
bridges, the geometric arctangent branch, and the optional Mathlib comparison.
The geometric $\pi$ node belongs to the equivalence statement and its proof,
not to the numerical-definition path for $\pi_{\log}$. Declaration and LOC
comparisons must include those bridges and report reuse of earlier chapters.
